\documentclass[11pt]{article}
\usepackage[a4paper,margin=25mm]{geometry}
\usepackage{amsmath,amssymb,amsthm,booktabs}
\usepackage[colorlinks=true,urlcolor=blue]{hyperref}
\newcommand{\Q}{\mathbb Q}\newcommand{\F}{\mathbb F}
\newcommand{\Z}{\mathbb Z}\newcommand{\OO}{\mathcal O}
\newcommand{\Sel}{\operatorname{Sel}}
\newcommand{\Cl}{\operatorname{Cl}}
\newtheorem{proposition}{Proposition}
\newtheorem{lemma}[proposition]{Lemma}
\emergencystretch=2em
\title{Excluding additional points in the four even sieve classes\\
for \(X_{\rm ns}^+(121)\)}
\author{}\date{29 October 2026}
\begin{document}\maketitle
\noindent\textbf{Status and scope.}
The new calculation excludes additional points in the sieve classes
\(n_0=-2,0,2,4\). The argument below uses the previously reproduced finite
sieve and its cusp-descent classes as input. The higher-prime local tables
in that sieve have not all been regenerated. The new even-case Selmer
calculations have been certified without GRH. The three odd classes
\(n_0=-1,1,3\) remain unresolved; this is not a determination of all
rational points on \(X_{\rm ns}^+(121)\).

\section{Statement and notation}
Put
\[
 E:y^2+y=x^3-x^2-7x+10,\qquad G=(4,5),\qquad E(\Q)=\Z G,
\]
and let \(K=\Q(z)\), where
\(z^5+z^4-4z^3-3z^2+3z+1=0\). Put
\[
 s=(-z^4-z^3+2z^2+3z+3,-2z^4-4z^3+6z^2+11z+1),
 \quad \alpha=x(2s).
\]
Here \(s\) has order 11. The finite sieve leaves precisely the states
\(n\equiv n_0\pmod M\), \(-2\le n_0\le4\), with
\(v_2(M)=4\) and \(v_{11}(M)=2\). The exact integer \(M\) is in
\path{outputs/sieve_result.json}. In each state, the cusp descent gives
\begin{equation}\label{eq:necessary}
 x(nG-3s)=\alpha+\delta_{n_0}w^{11},\qquad
 \delta_{n_0}=x(n_0G-3s)-\alpha,\qquad w\in K.
\end{equation}
This is the input from the preceding two reports.

\begin{proposition}\label{prop:main}
Let \(n_0\in\{-2,0,2,4\}\). If \(n\equiv n_0\pmod M\) and
\eqref{eq:necessary} holds for some \(w\in K\), then \(n=n_0\).
Consequently, subject to the inherited sieve input, any additional rational
point on \(X_{\rm ns}^+(121)\) must lie in one of the three odd states.
\end{proposition}

Set
\[
 L=K(\theta),\qquad \theta^3-\theta^2-7\theta+41/4=0,
 \qquad F=L(\zeta_{11}).
\]
The degrees are 15 and 30. Write \(\beta^2=\alpha-\theta\), with
\(\beta\in L\), and \(a=\delta_{n_0}/\beta^2\). In an even state the
two-descent gives \(v\in L\) with
\[
 v^2=\alpha-\theta+\delta_{n_0}w^{11}.
\]
After \(X=aw\), \(Y=a^5v/\beta\), this becomes
\begin{equation}\label{eq:C}
 C:\quad Y^2=X^{11}+b^2,\qquad b=a^5.
\end{equation}
Let \(P_0\in C(L)\) be the known point corresponding to \(n_0\), so
\(X(P_0)=a\), and embed \(C\) in \(J=\operatorname{Jac}(C)\) by
\(i(P)=[P-P_0]\). Write \(\sigma(X,Y)=(\zeta_{11}X,Y)\) and
\(\lambda=1-\sigma\). The group \(J[\lambda](F)\) is cyclic of order
11, generated by \(T=[(0,b)-\infty]\).

\section{The certified Selmer calculation}
At each of the three primes \(v\mid11\) of \(F\), the completion is
\(\Q_{11}(\zeta_{11})\), with residue field \(\F_{11}\). Denote the
Kummer map by \(\mu_v:J(F_v)\to J(F_v)/\lambda J(F_v)\), realized by
the descent functions \(Y-b\) and \(Y+b\).
Let \(\mathcal S\) be the image of \(\Sel^\lambda(J/F)\) under the
joint localization at these three places. Below we define a formal
subgroup \(H_{1,v}\); put
\[
 U=\prod_{v\mid11}\mu_v(H_{1,v}),\qquad
 T_{\rm loc}=\prod_{v\mid11}\langle\mu_v(T)\rangle,
 \qquad T_{\rm glob}=\langle(\mu_v(T))_v\rangle.
\]
The computed dimensions are
\[
\begin{array}{c|rrrr}
 n_0&-2&0&2&4\\\hline
 \dim\Sel^\lambda(J/F)&5&5&4&4\\
 \dim\ker(\mathrm{loc}_{11})&0&0&0&0\\
 \dim(\mathcal S\cap U)&0&0&0&0\\
 \dim(\mathcal S\cap T_{\rm loc})&1&1&1&1\\
 \dim(\mathcal S\cap(U+T_{\rm loc}))&1&1&1&1
\end{array}
\]
More precisely, exact linear algebra verifies
\begin{equation}\label{eq:intersection}
 U\cap T_{\rm loc}=0,\qquad
 \mathcal S\cap(U+T_{\rm loc})=T_{\rm glob}.
\end{equation}
Every local \(\mu_v(T)\) is nonzero, and \(\dim U=\dim T_{\rm loc}=3\).
The matrices are recorded in \path{formal_*_data.m};
\path{audit_matrices.m} verifies the identities independently of the
Selmer driver.

Here is how the GRH qualification from the previous report was removed.
The class numbers of \(L\) and \(H=\Q(\zeta_{11})\) are proved to be
one. The extension \(F/K\) is Galois with group \(S_3\), and \(L,H\)
are fixed fields of subgroups of orders 2 and 3. Restriction and norm,
with degrees prime to 11, identify the corresponding invariant
11-primary class groups with those of these subfields. Every irreducible
\(\F_{11}[S_3]\)-module has invariants under at least one of these two
subgroups. Hence \(\Cl(F)[11]=0\).

There are four primes in \(S_F\), namely the primes above 2 and 11.
Therefore \(\dim F(S_F,11)=14+1+4=19\). The split descent algebra is
\(F\times F\), so its global Kummer space has dimension 38. Although
GRH was used to discover candidate generators, the code checks that
each candidate has valuations divisible by 11 outside \(S_F\), and
that the 38 candidates have independent joint local images. This
certifies the full global space without GRH. The local images at 11
are generated by actual local points until their known full dimension
is reached; the routine also checks precision bounds for the Kummer
coordinates. The other local calculations use Magma's descent routines.
Thus the displayed even-case Selmer data are unconditional.

In particular, the previous rank upper bounds are now unconditional:
\[
 \operatorname{rank}J_{-2}(L),\operatorname{rank}J_0(L)\le20,
 \qquad \operatorname{rank}J_2(L),\operatorname{rank}J_4(L)\le15.
\]
The exclusion argument does not require these bounds to be sharp or a
finite-index Mordell--Weil subgroup.

\section{The local formal calculation}
Fix \(v\mid11\), write \(k=F_v\), \(\pi=1-\zeta_{11}\), and normalize
\(v_k(\pi)=1\). Then \(v_k(11)=10\) and
\(11/\pi^{10}\equiv-1\pmod\pi\). In all four cases,
\[
 \bar a=8,\qquad \bar b^2=1,\qquad \overline{Y(P_0)}\ne0.
\]
Use the proper hyperelliptic model of \eqref{eq:C} over \(\OO_k\).
Its special fibre is geometrically integral, with one cusp, and its
relative identity Picard group is smooth of dimension five. Let
\(H_m\) be its formal subgroup of points congruent to zero modulo
\(\pi^m\). Its cotangent module has integral basis
\[
 \omega_j=\frac{X^{j-1}dX}{2Y},\qquad 1\le j\le5.
\]
The automorphism \(\sigma\) extends to the model and reduces to the
identity. Consequently the power series for \(\lambda\) in formal
coordinates have all coefficients divisible by \(\pi\). Their linear
terms have eigenvalues \(1-\zeta_{11}^j\), whose quotients by \(\pi\)
reduce to \(j\). The formal inverse function theorem therefore gives
\begin{equation}\label{eq:division}
 \lambda:H_m\xrightarrow{\sim}H_{m+1}\qquad(m\ge1).
\end{equation}
On leading terms, dividing by \(\lambda\) divides the \(j\)-th
cotangent-dual coordinate by \(j\), after the common factor \(\pi\)
has been removed.

To compute \(\mu_v(H_1)\), choose smooth sections with distinct
residual \(X\)-coordinates
\(0,2,3,4,8\); at \(X=0\) choose \(Y=-b\).
The sum of the five Abel maps is a system of formal coordinates, since
its differential is a Vandermonde matrix with unit determinant.
For a displacement \(X\mapsto X+\pi u\), direct expansion gives
\[
 \frac{Y_{\rm new}-b}{Y_{\rm old}-b}
 \equiv
 \begin{cases}
 1+\pi^{11}\bar u/(4\bar b^2)&\bar X=0,\\
 1&\bar X\in\{2,3,4,8\}
 \end{cases}\pmod{\pi^{12}}.
\]
Indeed, at a nonzero residue the two possible terms of valuation 11
in \((X+\pi u)^{11}-X^{11}\) cancel:
\( (11/\pi^{10})\bar X^{10}\bar u+\bar u^{11}=0\).
All elements of \(1+\pi^{12}\OO_k\) are eleventh powers, by logarithm
and exponential on \(1+\pi^2\OO_k\). The first line generates the
nonzero unramified Kummer line. Its nonvanishing is also verified by
the local Kummer maps in the computation. Thus \(\dim\mu_v(H_1)=1\).

Let
\[
 \ell_0(X)=\frac{(X-2)(X-3)(X-4)(X-8)}{(-2)(-3)(-4)(-8)}
 =1+2X+2X^2+X^3+9X^4\in\F_{11}[X].
\]
This is the linear functional extracting the displacement at the section
with residual coordinate zero, expressed in the differential basis.
For a point \(P\) near \(P_0\), suppose \(v_k(X(P)-X(P_0))=d\ge2\).
Then \(i(P)\in H_d\setminus H_{d+1}\). Put
\(Q=\lambda^{-(d-1)}i(P)\in H_1\), using the unique formal divisions
in \eqref{eq:division}. Up to a nonzero scalar, its Kummer image is
\begin{equation}\label{eq:leading}
 h_{d-1}(8),\qquad
 h_r(x)=\sum_{j=1}^5\ell_{0,j-1}\,x^{j-1}j^{-r},
\end{equation}
where \(\ell_{0,j-1}\) is the coefficient of \(X^{j-1}\) in \(\ell_0\).
This follows by applying the diagonal leading action of successive
formal divisions to the Abel tangent vector
\((1,8,8^2,8^3,8^4)/(2\overline{Y(P_0)})\).
All nonlinear terms disappear on the associated graded groups.
In particular,
\begin{equation}\label{eq:five}
 r\equiv1\pmod{10}\quad\Longrightarrow\quad h_r(8)=h_1(8)=5\ne0.
\end{equation}

\section{From the rational-image locus to the formal obstruction}
Write \(u=2-z\), a uniformizer of the completion of \(K\) at 11,
normalized by \(v_K(u)=1\). The substitutions
\(x'=(x-4)/u^2\), \(y'=(y+1/2)/u^3\) give the integral minimal model
\[
 (y')^2=(x')^3+\frac{11}{u^2}(x')^2+
 \frac{33}{u^4}x'+\frac{121}{4u^6}.
\]
Its discriminant has valuation three. The point \(2G\) is formal, with
parameter \(-x'/y'=4u\). If \(n\ne n_0\), let
\(k_0=v_{11}(n-n_0)\ge2\). Since \(n-n_0\) is even, the elliptic
formal logarithm, which is an isometry here, gives
\[
 v_K(t((n-n_0)G))=1+5k_0.
\]
For \(q_0=n_0G-3s\), the exact computations give
\[
 v_K(x(q_0)-4)=2,\quad v_K(y(q_0)+1/2)=3,
 \quad v_K(\delta_{n_0})=2.
\]
Thus \(q_0\) is a smooth point on the special fibre of the minimal
model, with unit \(y'\), and translation by the formal point gives
\[
 v_K(x(nG-3s)-x(q_0))=3+5k_0.
\]
Equation \eqref{eq:necessary} now yields
\[
 v_K(w^{11}-1)=1+5k_0,
 \qquad v_K(w-1)=5k_0-4.
\]
For the second equality, first reduce modulo \(u\) to get \(\bar w=1\).
In the binomial expansion the term \(11(w-1)\) then has strictly
smaller valuation than all other nonzero terms, since \(v_K(11)=5\).
At every \(v\mid11\) of \(F\) we obtain
\begin{equation}\label{eq:d}
 d=v_k(X(P)-X(P_0))=10k_0-8\ge12,
 \qquad d-1\equiv1\pmod{10}.
\end{equation}

One must also check that the same global choice of square root puts
\(P\) near \(P_0\) at all three places. This is not automatic from
\(w\equiv1\). There is an explicit choice. Put
\[
 R=(n/2)G+4s,\qquad R_0=(n_0/2)G+4s,
\]
so \(2R=nG-3s\). Then
\begin{equation}\label{eq:root}
 v(R)=\frac{(x(R)-\theta)^2-(3\theta^2-2\theta-7)}{2y(R)+1}
\end{equation}
satisfies \(v(R)^2=x(2R)-\theta\). At \(R_0\) this equals the chosen
known square root exactly. Since \(4\mid M\), the difference
\(R-R_0=((n-n_0)/2)G\) is formal; its parameter has valuation at least
22 over each \(F_v\). All \(R_0\) have integral minimal coordinates.
The computed valuations are
\[
 4\le v_k(x(R_0)-\theta)\le6,\quad
 v_k(2y(R_0)+1)\le8,\quad v_k(v(R_0))=2.
\]
Translation in the minimal coordinates changes the original \(x\)
and \(y\) coordinates by quantities of valuations at least 26 and 28.
Substituting in \eqref{eq:root} shows
\(v(R)/v(R_0)\equiv1\pmod{\pi^{20}}\).
This fixes the required signs simultaneously. Since \(2Y(P_0)\) is
a unit, \(P\) is in its smooth residue disc, and \eqref{eq:d} applies.

\begin{lemma}\label{lem:global}
Under the Selmer injectivity and intersection conditions in
\eqref{eq:intersection}, no global divisor \(D\in J(F)\) can satisfy
\(D\in\prod_{v\mid11}H_{d,v}\), with \(d\ge2\), and have a nonzero
Kummer image after its \(d-1\) successive formal \(\lambda\)-divisions.
\end{lemma}
\begin{proof}
Write \(Q_j\) for the formal quotient \(\lambda^{-j}D\), for
\(0\le j\le d-1\). Initially \(Q_0=D\) is global, and its local
Kummer image is zero because \(H_d\subset\lambda H_1\).
Injectivity of Selmer localization implies that its global Kummer
class is zero. Hence a global \(R\) with \(\lambda R=Q_0\) exists.
Locally, \(R=Q_1+t_vT\), for scalars \(t_v\in\F_{11}\). Its local
Kummer image lies both in \(\mathcal S\) and in \(U+T_{\rm loc}\).
By \eqref{eq:intersection}, the formal contribution is zero and all
\(t_v\) coincide. Subtracting this global multiple of \(T\) shows
that \(Q_1\) is global and has zero global Kummer class, the latter
again by localization injectivity. Repeating the argument shows that
the local Kummer image of \(Q_{d-1}\) is zero, as asserted.
\end{proof}

\begin{proof}[Proof of Proposition~\ref{prop:main}]
If \(n\ne n_0\), the construction above gives \(D=i(P)\in J(F)\)
with depth \(d=10k_0-8\) at every place above 11.
Equations \eqref{eq:leading}--\eqref{eq:five} show that its final formal
quotient has nonzero Kummer image at every such place. This contradicts
Lemma~\ref{lem:global}.
\end{proof}

The global argument is a direct isogeny version of the division-image
method in Stoll's Selmer-group Chabauty framework \cite[Section~2]{stoll}.
The local formal calculation and the numerical certificates used here
are given explicitly above, rather than inferred from a rank inequality.

\section{The remaining odd cases and reproducibility}
For an odd state the same genus-five descent has
\((4-\theta)v^2=\alpha-\theta+\delta_{n_0}w^{11}\).
Its branch field is \(B=F(\sqrt{4-\theta})\), of degree 60.
The extension \(B/L\) is biquadratic, with degree-30 subfields
\[
 F=L(\sqrt{-11}),\quad L(\sqrt{4-\theta}),\quad
 L(\sqrt{-11(4-\theta)}).
\]
Since 11 does not divide four, the Kummer space of \(B\), unramified
outside 2 and 11, is generated by the images of the corresponding
spaces of these three subfields. This follows either from the four
characters of the Klein group or from
\(2=\sum_{|H|=2}N_H-N_{\operatorname{Gal}(B/L)}\) in its group ring.
The two new fields have computed class number one under GRH and
signatures \((10,10)\), \((0,15)\); their absolute discriminants are
\(2^{30}11^{28}\), \(2^{30}11^{27}\). Their unit computations have
completed. This supplies a way around the stalled direct degree-60
calculation. The corresponding subfield Kummer spaces have dimensions
19, 22 and 17. Their images should span a space of dimension 34 over
\(B\); the attempted combined localization has not yet certified that
span, so no odd Selmer bound is reported.

There is a concrete local difference in the odd cases. After adjoining
\(\sqrt{4-\theta}\), the unit model is again
\(Y^2=X^{11}+a^{10}\), but its known point reduces to
\((X,Y)=(-1,0)\), the singular point, at all three places above 11.
The extension \(B/F\) is ramified quadratic at two places and unramified
quadratic at the third; the normalized \(B\)-valuation of \(Y(P_0)\)
is one at all three. These facts are checked in
\path{verify_odd_geometry.m}. Thus the smooth-disc calculation used
above cannot simply be reused. A next local calculation should treat
the tangent direction at these singular reductions and compare its
division images at the two ramified places with the odd Selmer image.
No odd-class exclusion is claimed here.

The main scripts, run from the workspace root, are:
\begin{verbatim}
python3 outputs/continuation/create_descent.py
magma -b NN:=0 TAG:=0 outputs/continuation/compute_selmer.m \
  outputs/continuation/local_formal_test.m
magma -b outputs/continuation/formal_0_data.m \
  outputs/continuation/audit_matrices.m
magma -b outputs/continuation/verify_even_geometry.m
\end{verbatim}
Repeat the descent for \(n_0=-2,2,4\), using tags \texttt{minus2},
\texttt{2}, \texttt{4}. The four completed logs are
\path{selmer_formal_*.log}. The exact matrix audit is fast and does
not rerun number-field arithmetic. The generator script uses the user's
licensed Magma package to make a local copy with additional sampling
and certificate hooks; the installed package is untouched, and its
source is not redistributed. Magma V2.30-1 was used.

The finite sieve input retains the earlier audit scope: its linear algebra
and congruences were independently reproduced, and all 21 local tables
at eligible primes at most 200 were regenerated. The higher-prime tables
remain inherited data. The proof of Proposition~\ref{prop:main} itself
uses only the indicated congruence, descent equation, and the new
Selmer and local computations.

\begin{thebibliography}{1}
\bibitem{stoll} M. Stoll, \emph{Chabauty without the Mordell--Weil group},
preprint, December 2017.
\url{https://www.mathe2.uni-bayreuth.de/stoll/papers/ratpts-selmer-2016-12-07.pdf}.
\end{thebibliography}
\end{document}
